(1) 扰动满足
\[
\partial_t\mathbf u+2\boldsymbol\Omega\times\mathbf u=-\frac1\rho\nabla\pi,
\qquad\nabla\cdot\mathbf u=0.
\]
对 \(\mathbf u=\widehat{\mathbf u}e^{i(\mathbf k\cdot\mathbf r-\omega t)}\) 投影到垂直 \(\mathbf k\) 的平面，压力项消失。对 \(\boldsymbol\Omega=\Omega\mathbf e_z\)，旋转投影在该二维平面产生以 \(2\Omega\cos\theta\) 为系数的反对称矩阵，因此
\[
\boxed{\omega_\pm=\pm2\Omega\cos\theta}.
\]
这是带符号波动本征值；固定正频率可通过选择相应正负螺旋分支和 \(\mathbf k\) 方向获得。其绝对值至多 \(2\Omega\)，不同于随波数线性增长的可压缩声波。

(2) 在给定正交基中，\(P_\perp(\mathbf e_z\times\mathbf e_1)=\cos\theta\,\mathbf e_2\)，
\(P_\perp(\mathbf e_z\times\mathbf e_2)=-\cos\theta\,\mathbf e_1\)。
若 \(\cos\theta>0\)，正频率 \(\omega=2\Omega\cos\theta\) 的复振幅可取
\[
\widehat{\mathbf u}=A(\mathbf e_1-i\mathbf e_2).
\]
其横向两分量正交且相差四分之一周期，故轨迹为横向圆偏振。将波动方程沿 \(\mathbf k\) 取内积，得到
\[
\widehat\pi=\frac{2i\rho\Omega}{k^2}\,
\mathbf k\cdot(\mathbf e_z\times\widehat{\mathbf u})
=-\frac{2\rho\Omega A\sin\theta}{k}.
\]
可选择实 \(A\) 作相位规范。 \(\theta=0\) 时波矢沿转轴，频率为 \(2\Omega\)、压强扰动为0；\(\theta=\pi/2\) 时两分支合并为零频的地转退化模，应区分静态模式和非零频传播波。

(3) 取正向传播分支 \(\omega=2\Omega k_z/|\mathbf k|\)，则
\[
\mathbf c_{\rm ph}=\frac{\omega}{k}\widehat{\mathbf k},
\qquad
\mathbf c_g=\nabla_{\mathbf k}\omega
=\frac{2\Omega}{k}\left[\mathbf e_z-\cos\theta\,\widehat{\mathbf k}\right].
\]
直接相乘有
\(\mathbf k\cdot\mathbf c_g=2\Omega[\mathbf k\cdot\mathbf e_z-\cos\theta\,k]/k=0\)。
群速大小 \(|\mathbf c_g|=2\Omega|\sin\theta|/k\)；相速大小为 \(2\Omega|\cos\theta|/k\)。相速和群速虽都按 \(k^{-1}\) 缩放，方向却严格正交；这一结论源自频率仅依赖波矢方向。在有限容器、高波数黏性耗散、层化或非均匀转速环境中，频散需修正。

(4) 将Euler线性方程与 \(\rho\mathbf u\) 点乘，科里奥利力因
\(\mathbf u\cdot(\boldsymbol\Omega\times\mathbf u)=0\) 不做功，利用不可压性得到
\[
\partial_t\left(\frac{\rho|\mathbf u|^2}{2}\right)
+\nabla\cdot(\pi\mathbf u)=0.
\]
第二问的复模幅度为 \(A(\mathbf e_1-i\mathbf e_2)\)，因此
\[
\overline{\mathcal E}=\frac{\rho}{4}|\widehat{\mathbf u}|^2
=\frac{\rho|A|^2}{2}.
\]
又 \(\widehat\pi=-2\rho\Omega A\sin\theta/k\)，所以
\[
\overline{\mathbf S}=\frac12\Re(\widehat\pi\,\widehat{\mathbf u}^*)
=-\frac{\rho\Omega|A|^2\sin\theta}{k}\,\mathbf e_1.
\]
从第三问
\(\mathbf c_g=(2\Omega/k)(-\sin\theta)\mathbf e_1\)，恰有
\(\overline{\mathbf S}=\overline{\mathcal E}\mathbf c_g\)。
这也验证压力虽无净压缩功，却通过压力速度相关承担扰动能量输运。
